\section{Introduction}
There is significant need for flexible probability distributions in statistical domains. Although most frequently used densities are unimodal and ``bell shaped", their shapes can be arbitrarily complex. Several new candidates, such as \cite{chalabi2012flexible}, and at the same time, multiple new techniques, some of which are summarized in \cite{anusha2024dists}, have been devised to generate more flexible shapes from existing distributions.

One instinctive choice for a bell-shape is the crest of a sine-curve. Even disregarding its elegant mathematical properties, this crest forms a visually natural candidate as a probability density function. \cite{gilbert1893moon} mentions the \sindist{} in its most well recognized form to study the distribution of moon-craters, although in \cite{edwards2000sine} we find a brief history of the same density expression in the works of \cite{lagrange1770memoire} and \cite{pearson1902systematic}.

Several attempts have been made to extend the Sine Distribution into a more flexible array of shapes to model a wide variety of data distributions, including SS transformed Sine distribution (\cite{kumar2015bladder}), Sine-G distribution (\cite{mahmood2019sineg}), and Sine generalized family (\cite{oramulu2024sinegen}) although such recent modifications propose formulatically complex extensions of Gilbert's proposed elegant density. In our work, we consider a simple 4-parameter extension of Gilbert's baseline Sine Distribution into a flexible, location-scale-skewness-kurtosis family, that is equally elegant in its density expression to Gilbert's Sine Distribution.

We navigate through the work as follows: at first, we introduce the \sindist{} and investigate its shape properties. We will then take a look at some immediate mathematical results that follow from its density expression. Then we take a deeper look at its mathematical properties using the notion of Generalized Hypergeometric Functions. Then, we provide two instances of the flexibility of the \sindist{} by highlighting its fitting capabilities. Finally, we discuss its possible future applications.